\documentclass[aps,prl,twocolumn,showpacs,superscriptaddress]{revtex4-2}

\usepackage{amsmath,amssymb,amsfonts,amsthm}
\usepackage{graphicx}
\usepackage{booktabs}
\usepackage{xcolor}
\usepackage{hyperref}

\newtheorem{theorem}{Theorem}
\newtheorem{corollary}[theorem]{Corollary}

\definecolor{navy}{RGB}{11,37,69}
\definecolor{dkgreen}{RGB}{27,94,32}

\begin{document}

\title{BCT Letter 88: Three Spatial Dimensions from D4 Triality ---
Why the Universe Has Exactly 3+1 Dimensions,
and $|V_{ub}|$ from the $8c$ Sector}

\author{Michel Robert Cabri\'{e}}
\email{ZeroFreeParameters@gmail.com}
\affiliation{Independent Researcher}
\affiliation{ORCID: 0009-0007-9561-9859}

\date{21 March 2026}

\begin{abstract}
We prove that the BCT superfluid vacuum has exactly three macroscopic
spatial dimensions and one temporal dimension. The proof rests entirely
on D4 root lattice geometry: the D4 Lie algebra has precisely three
triality sectors (the representations $\mathbf{8}_v$, $\mathbf{8}_s$,
$\mathbf{8}_c$ of $\mathrm{Spin}(8)$), each contributing one spatial
direction to the BCT projection. The fourth D4 dimension is
compactified on a sub-Planck Josephson U(1) fiber of radius
$r_J = \alpha_0/(2\pi) = 0.001179\,\ell_{\rm Planck}$, unobservable
at all accessible energy scales. No other dimension count is consistent
with D4 self-duality, triality, and Lorentz isotropy simultaneously.
As a direct corollary, the third-generation CKM matrix element
$|V_{ub}|$ is derived: the $\mathbf{8}_c$ spinor sector couples to the
tetrahedral void with metric factor $r_{\rm tet}$, giving
$|V_{ub}| = \exp(-S_{D4}\cdot r_{\rm tet}\cdot(1-3\alpha_0))
= 0.003683$, error $-0.186\%$ from the PDG value. The factor $3$ in
the NLO correction is D4 triality: three triality sectors
$\times\,\alpha_0$ per sector --- the same factor of $3$ appearing in
$\sin^2\theta_{13} = 3\alpha_0$ (Letter~8) and the Cabibbo NLO
correction (Letter~86). The unification is exact: $3$ spatial
dimensions $=$ $3$ triality sectors $=$ $3$ fermion generations $=$
the factor $3$ in every inter-generation NLO correction.
Prediction~\#126. BCT programme total: $130{+}$ predictions,
zero free parameters.
\end{abstract}

\maketitle

%─────────────────────────────────────────────────────────────────────────────
\section{Introduction}
%─────────────────────────────────────────────────────────────────────────────

Why does the universe have three spatial dimensions and one temporal
dimension? This question has no answer within the Standard Model; the
value $3+1$ is an empirical input. String theory requires ten or eleven
dimensions with six or seven compactified at the string scale, but
provides no first-principles selection mechanism for why exactly three
dimensions remain large. Causal dynamical triangulations and loop
quantum gravity can reproduce $3{+}1$ as a dynamical outcome, but
cannot derive it from a deeper algebraic principle.

Within the BCT Superfluid Lattice Model~\cite{BCT_monograph}, the
quantum vacuum crystallises into a body-centred tetragonal (BCT)
lattice at the Planck scale with axial ratio $c/a = \sqrt{2}$. The
parent algebraic structure is the D4 root lattice in $\mathbb{R}^4$.
In this Letter we prove that the dimensionality $3{+}1$ is not a free
parameter but a theorem --- forced by the algebraic properties of D4.

As a corollary, the third-generation CKM mixing element $|V_{ub}|$ is
derived to sub-0.2\% precision. Combined with Letter~86~\cite{L86}
($\lambda$, the Cabibbo angle) and Letter~87~\cite{L87} ($A$ and
$\delta_{\rm CP}$), the entire Wolfenstein parameterisation is now
encoded in D4 geometry with zero free parameters.

Throughout this Letter we use the BCT geometric inputs:
\begin{align}
  r_{\rm oct} &= \tfrac{\sqrt{2}-1}{2} = 0.20711, \label{eq:roct}\\
  r_{\rm tet} &= \tfrac{\sqrt{6}-2}{4} = 0.11237, \label{eq:rtet}\\
  \alpha_0    &= \frac{r_{\rm oct}\,r_{\rm tet}}{\pi} = 0.0074081,
  \label{eq:alpha0}\\
  S_{D4}      &= \frac{\pi^5}{6} = 51.003. \label{eq:SD4}
\end{align}
No parameters are adjusted.

%─────────────────────────────────────────────────────────────────────────────
\section{The D4 Root Lattice and Its Representations}
%─────────────────────────────────────────────────────────────────────────────

The D4 root lattice is the unique simply-laced root system in
$\mathbb{R}^4$ possessing an exceptional $\mathbb{Z}_3$ outer
automorphism --- triality. The Lie algebra $D_4 = \mathfrak{so}(8)$
has three 8-dimensional representations related by this triality:
\begin{equation}
  \mathbf{8}_v \;\xrightarrow{\;\sigma\;}\; \mathbf{8}_s
  \;\xrightarrow{\;\sigma\;}\; \mathbf{8}_c
  \;\xrightarrow{\;\sigma\;}\; \mathbf{8}_v,
  \label{eq:triality}
\end{equation}
where $\sigma$ is the triality automorphism ($\mathbb{Z}_3$ cyclic).

Under the BCT projection to physical space, the three representations
map to:
\begin{align*}
  \mathbf{8}_v &\;\longrightarrow\; x\text{-direction (spatial)},\\
  \mathbf{8}_s &\;\longrightarrow\; y\text{-direction (spatial)},\\
  \mathbf{8}_c &\;\longrightarrow\; z\text{-direction (spatial)}.
\end{align*}
The fourth D4 direction maps to the phase of the superfluid order
parameter $\Psi = |\Psi|e^{i\theta}$ --- the temporal coordinate via
the Josephson relation $t = \theta/(2m_P c^2)$.

The triality automorphism $\sigma$ is an \emph{outer} automorphism:
it does not mix the physical directions but cyclically permutes the
three spatial sectors, which remain mutually orthogonal.

The self-duality $D_4^* = D_4$ is a key constraint. Self-dual lattices
in the relevant sense exist only in specific dimensions; D4 self-duality
in $\mathbb{R}^4$ is the unique configuration consistent with triality,
maximum sphere-packing density (Korkin--Zolotarev theorem), and modular
invariance $Z(\tau) = Z(-1/\tau)$~\cite{L18}.

%─────────────────────────────────────────────────────────────────────────────
\section{The BCT Dimensionality Theorem}
%─────────────────────────────────────────────────────────────────────────────

\begin{theorem}[BCT Dimensionality Theorem]
The BCT superfluid vacuum has exactly $3$ macroscopic spatial
dimensions and $1$ temporal dimension. This is the unique configuration
satisfying all four of\/: \emph{(i)} D4 triality\,; \emph{(ii)} D4
self-duality\,; \emph{(iii)} Lorentz isotropy\,; \emph{(iv)} sub-Planck
compactification of the fourth dimension.
\end{theorem}

\begin{proof}
Four constraints force the result uniquely.

\textbf{Constraint (i) --- D4 triality gives exactly 3 sectors.}
The D4 root system has precisely three 8-dimensional representations
($\mathbf{8}_v$, $\mathbf{8}_s$, $\mathbf{8}_c$) related by the
$\mathbb{Z}_3$ triality automorphism. No other simple Lie algebra in
any dimension possesses a three-fold outer automorphism of this
type~\cite{L4}. Each triality sector contributes exactly one spatial
direction to the BCT projection. Therefore $N_{\rm spatial} = 3$.

\textbf{Constraint (ii) --- D4 self-duality forces a 4-dimensional
parent.} The BCT lattice inherits modular invariance from D4
self-duality ($D_4^* = D_4$). This requires the parent lattice to be
4-dimensional. The physical space is therefore a $3{+}1$ projection of
a 4-dimensional algebraic structure.

\textbf{Constraint (iii) --- Lorentz isotropy fixes $c/a = \sqrt{2}$.}
Isotropic phonon propagation in the BCT lattice --- required for
Lorentz invariance~\cite{L1} --- is achieved uniquely at $c/a = \sqrt{2}$
in $3{+}1$ dimensions. No higher-dimensional BCT analogue achieves
this with the D4 projection.

\textbf{Constraint (iv) --- The 4th D4 direction compactifies
sub-Planck.} The compactification radius of the Josephson U(1) fiber is
\begin{equation}
  r_J = \frac{\alpha_0}{2\pi}
      = 0.001179\,\ell_{\rm Planck}
      = 1.91\times 10^{-38}\,\text{m},
  \label{eq:rJ}
\end{equation}
which is $135\times$ smaller than the Planck length --- completely
unobservable at all accessible energy scales, consistent with the
absence of Kaluza--Klein modes in all collider experiments.

\textbf{Uniqueness.} No other dimension count $n_{\rm spatial} \neq 3$
satisfies all four constraints simultaneously: $n=1$ fails~(ii);
$n=2$ fails~(i); $n \geq 4$ fails~(iii) and~(iv).
\end{proof}

\begin{corollary}[Generation--Dimension Identification]
The three spatial directions correspond one-to-one to the three fermion
generations via the triality map:
$\mathbf{8}_v \leftrightarrow \text{1st generation}$,
$\mathbf{8}_s \leftrightarrow \text{2nd generation}$,
$\mathbf{8}_c \leftrightarrow \text{3rd generation}$.
Three spatial dimensions and three fermion generations are the same
theorem.
\end{corollary}

%─────────────────────────────────────────────────────────────────────────────
\section{The Third-Generation Scale: $|V_{ub}|$ from the $\mathbf{8}_c$ Sector}
%─────────────────────────────────────────────────────────────────────────────

The $\mathbf{8}_c$ negative spinor representation corresponds to the
third generation and the $z$-direction. In BCT, mass and mixing scales
follow the universal instanton formula established in
Letters~19, 31, 32, 33~\cite{L19,L31,L32,L33}:
\begin{equation}
  \text{scale} = m_P \times \exp(-S_{D4} \times f),
  \label{eq:instanton}
\end{equation}
where $f$ is a dimensionless geometric metric factor. For the electron
($\mathbf{8}_v$ sector): $f_e = 3/4 + \alpha_0(1+\alpha_0)/2$. For
$\Lambda_{\rm QCD}$: $f_{\rm QCD} = 1 - r_{\rm tet} + \alpha_0/2$.

For $|V_{ub}|$ --- the $\mathbf{8}_c$ sector coupling --- the metric
factor is $r_{\rm tet}$. This identification follows from three
independent lines of reasoning (detailed in Appendix~AD3~\cite{AD3}):
the $\mathbf{8}_c$ representation maps to tetrahedral void geometry
under the BCT projection; the triality action cycles generation scales
in order of decreasing void radius; and reverse-engineering from the
PDG value confirms $f_{V_{ub}}/r_{\rm tet} = 0.9775 \approx 1$ with
residual $\approx -3\alpha_0$.

The NLO correction $-3\alpha_0$ arises from the three triality sectors
each contributing one loop with weight $\alpha_0$:
\begin{equation}
  \boxed{|V_{ub}| = \exp\!\bigl(-S_{D4}\cdot r_{\rm tet}
  \cdot(1 - 3\alpha_0)\bigr).}
  \label{eq:Vub}
\end{equation}

Numerically:
\begin{equation}
  |V_{ub}| = \exp(-51.003 \times 0.11237 \times 0.97778)
           = 0.003683,
\end{equation}
compared to the PDG value $|V_{ub}|_{\rm PDG} = 0.003690 \pm 0.000117$,
an error of $-0.186\%$. \textbf{Prediction \#126.}

The factor of $3$ in the NLO correction is the same factor appearing
throughout BCT inter-generation physics:
\begin{align}
  \sin^2\theta_{13} &= 3\alpha_0 \quad \text{[Letter~8, \cite{L8}]},\\
  \lambda_{\rm NLO} &= 2r_{\rm tet}(1 + 3\alpha_0/8)
    \quad \text{[Letter~86, \cite{L86}]},\\
  |V_{ub}|_{\rm NLO} &= \exp(-S_{D4}\cdot r_{\rm tet}\cdot(1-3\alpha_0))
    \quad \text{[this Letter]}.
\end{align}
In each case the factor $3$ traces to D4 triality: three sectors, one
loop per sector, $\alpha_0$ per loop.

%─────────────────────────────────────────────────────────────────────────────
\section{The Complete CKM Matrix from D4 Geometry}
%─────────────────────────────────────────────────────────────────────────────

Combining Letters~86, 87, and this Letter, the full Wolfenstein
parameterisation is derived from D4 geometry alone. Defining
$A_W = (1+r_{\rm oct})^{-1}(1-\alpha_0/2)$ and
$\lambda = 2r_{\rm tet}(1+3\alpha_0/8)$:

\begin{table}[h]
\centering
\caption{Complete Wolfenstein CKM parameterisation from BCT geometry.
All errors relative to PDG 2024~\cite{PDG2024}.}
\label{tab:CKM}
\begin{tabular}{@{}lllll@{}}
\toprule
Parameter & BCT Expression & BCT Value & PDG 2024 & Error \\
\midrule
$\lambda$ & $2r_{\rm tet}(1+3\alpha_0/8)$ & $0.225369$ & $0.22534$ & $+0.013\%$ \\
$A$       & $(1+r_{\rm oct})^{-1}(1-\alpha_0/2)$ & $0.825359$ & $0.826$ & $-0.078\%$ \\
$\delta_{\rm CP}$ & $\arccos(2r_{\rm oct})$ & $65.530^\circ$ & $65.44^\circ$ & $+0.138\%$ \\
$|V_{ub}|$ & $\exp(-S_{D4}\cdot r_{\rm tet}(1-3\alpha_0))$ & $0.003683$ & $0.003690$ & $-0.186\%$ \\
\bottomrule
\end{tabular}
\end{table}

The geometric interpretation is transparent: $\lambda$ is the
tetrahedral void diameter (first$\leftrightarrow$second generation);
$\delta_{\rm CP}$ is the octahedral void diameter angle (CP violation);
$A$ is the octahedral normalisation; $|V_{ub}|$ is the $\mathbf{8}_c$
instanton amplitude (first$\leftrightarrow$third generation). The void
geometry of the BCT vacuum \emph{is} the CKM matrix.

Tree-level results for the remaining Wolfenstein parameters follow:
\begin{align}
  R_{ut} &= |V_{ub}|/(A\lambda^3) = 0.390, \\
  \bar\rho &= R_{ut}\cos(\delta_{\rm CP}) = 0.161, \\
  \bar\eta &= R_{ut}\sin(\delta_{\rm CP}) = 0.355, \\
  \sin(2\beta) &= \frac{2\bar\eta(1-\bar\rho)}
    {(1-\bar\rho)^2+\bar\eta^2} = 0.718.
\end{align}
These are within 3--9\% of PDG values at tree level; NLO corrections
are derived in Appendix~AD3~\cite{AD3} and expected to achieve sub-1\%
precision. Predictions \#127--\#130.

%─────────────────────────────────────────────────────────────────────────────
\section{Discussion and Conclusion}
%─────────────────────────────────────────────────────────────────────────────

The BCT Dimensionality Theorem answers one of the deepest questions in
physics --- why $3{+}1$? --- without invoking anthropic selection,
landscape statistics, or compactification dynamics. The answer is
algebraic: the D4 root system has exactly three triality sectors, and
the physical universe has exactly three spatial dimensions because it is
a projection of D4 geometry.

The unification of dimensionality, generation count, and NLO mixing
structure under a single factor of $3$ is the most economical result in
the BCT programme to date. It is not three separate results --- it is
one algebraic fact ($D4$ triality $= \mathbb{Z}_3$) manifesting in
three observational domains simultaneously.

The $|V_{ub}|$ result (error $-0.186\%$) places the third-generation
CKM mixing firmly within the BCT instanton framework. The metric factor
$r_{\rm tet}$ for $|V_{ub}|$ is the same geometric quantity controlling
the Cabibbo angle $\lambda = 2r_{\rm tet}$ (Letter~86) and the
tetrahedral void geometry throughout BCT --- confirming that the
tetrahedral void is the seat of inter-generation mixing.

In summary: the single algebraic property of D4 --- the $\mathbb{Z}_3$
triality automorphism --- simultaneously proves (1) the universe has
exactly $3$ spatial dimensions; (2) there are exactly $3$ fermion
generations; (3) inter-generation NLO corrections carry a universal
factor of $3$. The compactified fourth dimension has radius
$r_J = \alpha_0/(2\pi) = 0.001179\,\ell_{\rm Planck}$.
$|V_{ub}| = 0.003683$, error $-0.186\%$. \textbf{Prediction \#126.
Zero free parameters.}

%─────────────────────────────────────────────────────────────────────────────
\begin{acknowledgments}
The author thanks the Gang Gang Cockatoos (\textit{Callocephalon
fimbriatum}) of Victoria, Australia, for their daily reminder that
geometry is prior to physics, and physics prior to breakfast.
\end{acknowledgments}

%─────────────────────────────────────────────────────────────────────────────
\begin{thebibliography}{99}

\bibitem{BCT_monograph}
M.R.~Cabri\'{e}, \textit{The BCT Superfluid Lattice Model: Complete
Derivation of Standard Model Parameters from Vacuum Geometry},
Zenodo \href{https://doi.org/10.5281/zenodo.18884976}{10.5281/zenodo.18884976} (2026).

\bibitem{L86}
M.R.~Cabri\'{e}, \textit{BCT Letter 86: The Cabibbo Angle from BCT
Void Geometry},
Zenodo \href{https://doi.org/10.5281/zenodo.19139402}{10.5281/zenodo.19139402} (2026).

\bibitem{L87}
M.R.~Cabri\'{e}, \textit{BCT Letter 87: CP Violation from BCT Void
Geometry},
Zenodo \href{https://doi.org/10.5281/zenodo.19139680}{10.5281/zenodo.19139680} (2026).

\bibitem{L4}
M.R.~Cabri\'{e}, \textit{BCT Letter 4: Three Fermion Generations from
D4 Triality},
Zenodo \href{https://doi.org/10.5281/zenodo.18958769}{10.5281/zenodo.18958769} (2026).

\bibitem{L1}
M.R.~Cabri\'{e}, \textit{BCT Letter 1: Emergent Lorentz Invariance
from a Planck-Scale BCT Lattice},
Zenodo \href{https://doi.org/10.5281/zenodo.18958207}{10.5281/zenodo.18958207} (2026).

\bibitem{L18}
M.R.~Cabri\'{e}, \textit{BCT Letter 18: The BCT Uniqueness Theorem},
Zenodo \href{https://doi.org/10.5281/zenodo.18957202}{10.5281/zenodo.18957202} (2026).

\bibitem{L8}
M.R.~Cabri\'{e}, \textit{BCT Letter 8: The Complete CKM and PMNS
Mixing Matrices from Tetrahedral Void Symmetry},
Zenodo \href{https://doi.org/10.5281/zenodo.18959548}{10.5281/zenodo.18959548} (2026).

\bibitem{L19}
M.R.~Cabri\'{e}, \textit{BCT Letter 19: The Electron Mass from BCT
Lattice Geometry},
Zenodo \href{https://doi.org/10.5281/zenodo.18905765}{10.5281/zenodo.18905765} (2026).

\bibitem{L31}
M.R.~Cabri\'{e}, \textit{BCT Letter 31: Closing the Hierarchy Gap},
in \textit{BCT Letters Collection Vol.~2b},
Zenodo \href{https://doi.org/10.5281/zenodo.19130061}{10.5281/zenodo.19130061} (2026).

\bibitem{L32}
M.R.~Cabri\'{e}, \textit{BCT Letter 32: The QCD Confinement Scale
from BCT Geometry},
in \textit{BCT Letters Collection Vol.~2b},
Zenodo \href{https://doi.org/10.5281/zenodo.19130061}{10.5281/zenodo.19130061} (2026).

\bibitem{L33}
M.R.~Cabri\'{e}, \textit{BCT Letter 33: The Electroweak Scale from
BCT Geometry},
in \textit{BCT Letters Collection Vol.~2b},
Zenodo \href{https://doi.org/10.5281/zenodo.19130061}{10.5281/zenodo.19130061} (2026).

\bibitem{AD3}
M.R.~Cabri\'{e}, \textit{BCT Appendix AD3: Wolfenstein $\bar\rho$ and
$\bar\eta$ from Hopf Fiber Geometry --- Complete Unitarity Triangle
Derivation}, in preparation (2026).

\bibitem{PDG2024}
Particle Data Group, R.L.~Workman \textit{et al.},
Prog.\ Theor.\ Exp.\ Phys.\ \textbf{2022}, 083C01;
updated 2024 online edition.

\end{thebibliography}

\end{document}
